\section{Tema 1}
\begin{enumerate}[label=\arabic*.]

    \item \textbf{Crecimiento del PIB per cápita en etapas de expansión y crisis (España vs. Eurozona):} \\
    Durante el periodo 1961-2024, en las etapas de expansión económica, España suele registrar tasas anuales de variación del PIB per cápita real superiores a la media observada en la Eurozona (pág. 8). Por el contrario, al acontecer fases de crisis, la economía española tiende a sufrir los impactos con una intensidad significativamente mayor que sus socios comunitarios (pág. 17) , manifestando fluctuaciones del PIB real per cápita mucho más marcadas y agudas en todas las etapas recesivas (pág. 17).

    \item \textbf{La regla del 4\% sobre el saldo exterior:} \\
    Esta regla describe un comportamiento estructural observado tras la plena integración de España en la Unión Económica y Monetaria (UEM), según el cual la economía nacional adquirió la capacidad de sostener déficits comerciales elevados y persistentes, superiores al 4\% del PIB, sin verse obligada a recurrir a las devaluaciones cambiarias tradicionales (pág. 49). Este patrón fue viable gracias a la libre circulación de capitales que permitió captar de manera continuada la financiación necesaria en los mercados financieros internacionales (pág. 49).

    \item \textbf{Comportamiento del saldo exterior según el ciclo económico (1961-2024):} \\
    El saldo exterior español muestra una estrecha dependencia cíclica:
    \begin{itemize}
        \item En fases de fuerte expansión económica, el saldo se deteriora con rapidez hacia terrenos negativos (déficit abultado) debido a la intensa presión de la demanda interna y la consecuente inflación diferencial (pág. 13, 42).
        \item En etapas de crisis económica, el desplome de la demanda interna contrae severamente las importaciones, forzando de manera simultánea a las empresas a orientar su producción hacia mercados exteriores; esto corrige el desequilibrio hasta el punto de generar superávits comerciales continuados, tal como se constató a partir del año 2012 (pág. 13, 49).
    \end{itemize}

    \item \textbf{Incremento considerable del déficit exterior y sus razones:} \\
    Dentro del periodo histórico analizado, el déficit exterior experimentó un crecimiento sin precedentes durante la fase de expansión sostenida que culminó en el año 2007, momento en el cual el saldo negativo de la balanza por cuenta corriente alcanzó un histórico 10\% del PIB (pág. 42). Las razones fundamentales que explican este desequilibrio fueron:
    \begin{itemize}
        \item La extraordinaria fortaleza de la demanda interna, la cual sobreestimuló las importaciones en detrimento de las ventas al exterior (pág. 42).
        \item Una inflación diferencial persistente (situada en torno a 1,5 puntos por encima de la media europea durante casi todo el ciclo) que encareció los productos nacionales (pág. 42).
        \item Una productividad prácticamente estancada que impidió absorber los crecientes costes laborales unitarios, trasladándolos directamente a los precios y mermando la competitividad empresarial (pág. 42).
        \item La intensa presión competitiva ejercida en los mercados globales por el ascenso de las economías emergentes de Asia y Europa del Este (pág. 42).
    \end{itemize}

    \item \textbf{Disparidad del PIB per cápita frente a la Eurozona y la UE-27:} \\
    El PIB per cápita de España (en paridades de poder adquisitivo) se sitúa en una posición relativa inferior al medirse frente a la Eurozona (86,7\% en 2024) en comparación con su posición frente al conjunto de la UE-27 (90,0\% en 2024) (pág. 10). Esta diferencia se explica porque la Eurozona congrega de forma exclusiva a las economías con mayor nivel de renta y desarrollo del continente, mientras que la media de la UE-27 incluye a un grupo de países de incorporación más reciente con niveles de renta media notablemente inferiores, lo que eleva estadísticamente la posición relativa de España (pág. 10).

    \item \textbf{Existencia de convergencia real desde los años 60:} \\
    Sí, los datos históricos confirman un proceso sostenido de convergencia real a largo plazo (pág. 38). Partiendo de un PIB per cápita que apenas superaba el 70\% de la media continental a comienzos de los años sesenta, la economía española avanzó firmemente hasta alcanzar el 86,7\% de la Eurozona y el 90,0\% de la UE-27 en 2024 (pág. 10) , superando contratiempos y fases de alejamiento temporal provocados por las crisis estructurales de 1974, 2008 y 2020 (pág. 10).

    \item \textbf{Rasgos del crecimiento económico español:}
    \begin{itemize}
        \item Los países europeos han crecido a un ritmo mayor que España entre 1960 y 2024: \textbf{NO} (España promedió un crecimiento del 2,4\% anual acumulativo frente al 2,0\% de la Eurozona) (pág. 16).
        \item España ha participado intensamente en las relaciones económicas con los países de su entorno, incluso antes de pertenecer a la UE: \textbf{SÍ} (pág. 16).
        \item España destruye menos empleo que los países europeos durante las crisis: \textbf{NO} (debido a rigideces y a la alta temporalidad, el desempleo reacciona de forma exageradamente destructiva en las recesiones) (pág. 50, 60).
        \item España ha sufrido fluctuaciones más intensas del PIB que los países europeos en las etapas de crisis: \textbf{SÍ} (pág. 17).
    \end{itemize}

    \item \textbf{Evaluación del periodo desde los años 80 hasta el Tratado de Maastricht (1992):}
    \begin{itemize}
        \item A principios de los 80 tuvimos que acometer una reconversión industrial: \textbf{V} (pág. 26).
        \item Las empresas públicas estaban saneadas: \textbf{F} (constituían un grave problema estructural de ineficiencia) (pág. 26).
        \item Los servicios tenían una baja productividad: \textbf{V} (pág. 26).
        \item El mercado de trabajo funcionaba sin rigidez, permitiendo un ajuste perfecto entre la oferta y la demanda: \textbf{F} (la rigidez laboral era uno de los problemas de fondo fundamentales) (pág. 25, 26).
        \item El objetivo principal del gobierno socialista en 1982 era crear empleo, ya que estaba controlada la inflación, el déficit público y el déficit exterior: \textbf{F} (el fin último era crear empleo, pero la inflación y los déficits estaban descontrolados; reducirlos eran los objetivos intermedios obligatorios) (pág. 25).
        \item La Seguridad Social era deficitaria: \textbf{V} (pág. 25, 26).
        \item A partir de 1986 se inicia un periodo de recuperación económica, con unos desequilibrios macroeconómicos que se controlaron y redujeron en los cuatro años anteriores: \textbf{V} (gracias a las severas políticas de ajuste previas) (pág. 27).
        \item El periodo entre 1986 y 1991 se debió al empuje de la demanda interna (consumo + inversión) y de la demanda externa: \textbf{F} (se debió en exclusiva a la demanda interna; la externa restó crecimiento al generar un profundo déficit comercial) (pág. 28, 29).
        \item La inflación creció desde 1989 a 1991: \textbf{V} (los precios repuntaron hasta niveles cercanos al 7\%) (pág. 29, 30).
        \item La inflación de carácter dual se refiere a que los precios de los bienes crecen más que los precios de los servicios cuando hay un aumento de la demanda: \textbf{F} (significa lo contrario: los precios de los servicios crecen más rápido que los de los bienes comercializables) (pág. 29).
        \item La inflación de carácter diferencial se refiere a que los precios crecen más en España que en la media de países europeos cuando la economía crece: \textbf{V} (pág. 29).
        \item A partir de 1988 se aplicó una política monetaria restrictiva: \textbf{V} (pág. 29, 30).
        \item La cotización de la peseta aumentó tras el aumento de los tipos de interés a partir de 1988, aumentando nuestra competitividad frente al exterior: \textbf{F} (la fuerte apreciación de la peseta redujo drásticamente la competitividad exterior) (pág. 29, 30).
    \end{itemize}

    \item \textbf{Préstamos al sector público en la segunda fase de la UEM:} \\
    \textbf{NO} (pág. 31). Con el inicio de esta fase el 1 de enero de 1994, quedó prohibido de manera taxativa que los bancos centrales nacionales concedieran cualquier tipo de crédito o facilidades de financiación al sector público, una medida institucional estricta diseñada para forzar la disciplina fiscal y la consolidación de las cuentas públicas en los países miembros (pág. 31).

    \item \textbf{Motivo de exigencia de independencia de los bancos centrales en la 2ª fase:} \\
    La independencia se exigía como un requisito legal y operativo de carácter preparatorio (pág. 31). El objetivo era desvincular la política monetaria del control político directo, un proceso institucional obligatorio que debía completarse antes de la constitución definitiva del Sistema Europeo de Bancos Centrales (SEBC) para asegurar el correcto funcionamiento de la política monetaria única en la tercera fase (pág. 31).

    \item \textbf{Mecanismo de tipos de cambio europeo MTC II:} \\
    Es la estructura monetaria y cambiaria que entró en vigor el 1 de enero de 1999 coordinadamente con la tercera fase de la UEM (pág. 31). Su propósito formal es regular, supervisar y estabilizar los tipos de cambio de las divisas pertenecientes a los países de la Unión Europea que no han adoptado plenamente el euro, vinculándolas a la moneda única dentro de unas bandas de fluctuación definidas para evitar perturbaciones en el mercado común (pág. 31).

    \item \textbf{Pacto de Estabilidad y Crecimiento:} \\
    Es un acuerdo normativo e internacional adoptado formalmente al inicio de la tercera fase de la UEM (1999) (pág. 31). Establece un marco rígido de disciplina presupuestaria a nivel comunitario, fijando límites estrictos al déficit público y a la deuda soberana de los Estados miembros con el fin de evitar conductas fiscales expansivas que pongan en riesgo la estabilidad macroeconómica del euro (pág. 31, 42).

    \item \textbf{Criterios exigidos para acceder a la 3ª fase de la UEM:} \\
    Se exigía el estricto cumplimiento de los criterios de convergencia nominal y la correspondiente adecuación del marco institucional (pág. 33). Los requisitos nominales obligatorios eran:
    \begin{itemize}
        \item \textbf{Tasa de Inflación:} No superar en más de 1,5 puntos porcentuales la media de los tres países de la UE con mejor comportamiento de precios (pág. 33).
        \item \textbf{Tipos de Interés:} No superar en más de 2 puntos porcentuales el rendimiento de la deuda pública a largo plazo de los tres Estados más estables en precios (pág. 33).
        \item \textbf{Déficit Público:} No exceder bajo ningún concepto el 3\% del PIB (pág. 33).
        \item \textbf{Deuda Pública:} No superar el umbral de referencia fijado en el 60\% del PIB (pág. 33).
        \item \textbf{Estabilidad Cambiaria:} Mantener la moneda dentro de las bandas normales de fluctuación del SME sin tensiones ni devaluaciones (pág. 33).
    \end{itemize}

    \item \textbf{Cumplimiento de los requisitos por España al cierre de 1998:} \\
    \textbf{NO} de manera integral (pág. 37). Aunque España superó con éxito el examen oficial al cumplir los criterios de Inflación (1,4\%), Tipos de Interés (4,0\%) y Déficit Público (1,8\%), incumplió formalmente el límite rígido de la Deuda Pública, la cual se situó en un 67,4\% del PIB, superando el 60\% máximo establecido por el Tratado (pág. 37).

    \item \textbf{Razones del colapso del SME entre 1992 y 1993:} \\
    El colapso sistémico fue provocado por la conjunción de los severos desequilibrios macroeconómicos acumulados en las economías periféricas de Europa (pág. 35)  y las repercusiones financieras e internacionales derivadas de la reunificación de Alemania (pág. 35). Para financiar su proceso de unificación, las autoridades alemanas elevaron drásticamente sus tipos de interés, forzando una masiva e insostenible reconfiguración de los flujos de capital dentro del continente (pág. 34, 35).

    \item \textbf{Impacto de la elevación de tipos en Alemania sobre España (1992-1994):} \\
    La política de altos tipos de interés en Alemania provocó una desconfianza generalizada en los mercados y desencadenó una huida masiva de capitales fuera de España (pág. 35). A pesar de las intensas intervenciones del Banco de España, que quemó cuantiosas reservas vendiendo marcos y dólares para defender la moneda, la presión de los mercados resultó irresistible y abocó a una depreciación forzosa, concretada en tres devaluaciones sucesivas de la peseta que acumularon un recorte del 20\% de su valor entre 1992 y 1995 (pág. 35).

    \item \textbf{Hechos acontecidos en la crisis de 1992-1994:}
    \begin{itemize}
        \item Descenso del crecimiento del PIB sin llegar a ser negativo: \textbf{NO} (se registraron tasas reales negativas de variación de la producción) (pág. 36).
        \item Apreciación del tipo de cambio: \textbf{NO} (aconteció una severa depreciación y devaluación cambiaria) (pág. 35, 36).
        \item Caída de la FBK (inversión): \textbf{SÍ} (se experimentó un hundimiento muy acusado de la inversión en bienes de equipo y capital fijo) (pág. 36).
        \item Desempleo, inflación y déficit público: \textbf{SÍ} (el paro y la inflación superaron el 5\%, mientras que el déficit escaló por encima del 7\% del PIB) (pág. 36).
        \item Déficit exterior: \textbf{SÍ} (la balanza por cuenta corriente operó en situación de déficit permanente) (pág. 36).
    \end{itemize}

    \item \textbf{Cambio en los instrumentos de política económica desde el año 2000:} \\
    \textbf{SÍ}, se produjo una alteración estructural e irreversible de las herramientas tradicionales de gestión macroeconómica (pág. 39). Al formalizarse la zona Euro, las competencias sobre la política monetaria y la fijación del tipo de cambio fueron transferidas de manera exclusiva al Banco Central Europeo (BCE) (pág. 39). En consecuencia, los países miembros perdieron estas palancas y se vieron obligados a concentrar toda su capacidad económica en la consolidación fiscal rigurosa y en la ejecución de reformas estructurales internas (pág. 39).

    \item \textbf{Factores expansivos durante el periodo de crecimiento 1995-2007:}
    \begin{itemize}
        \item \textbf{Creación masiva de empleo:} Se incorporaron 7,84 millones de ocupados gracias a las reformas laborales de 1994 y 1996 y a una notable moderación de los costes laborales unitarios (pág. 40).
        \item \textbf{Incremento del consumo privado:} Impulsado por la mejora sostenida de la renta familiar disponible, el fuerte efecto riqueza patrimonial (burbuja inmobiliaria y bursátil) y la generalización de tipos de interés reales bajos con abundante liquidez internacional (pág. 41).
        \item \textbf{Aumento de la FBK Fijo:} Estimulado por la necesidad de dotaciones en capital público, el boom de la construcción residencial y la fuerte adquisición de bienes de equipo asociados a la expansión (pág. 41).
    \end{itemize}

    \item \textbf{Factores contractivos durante el periodo 1995-2007:}
    \begin{itemize}
        \item \textbf{Política Fiscal Restrictiva:} La obligación de dar cumplimiento al Pacto de Estabilidad y a la posterior Ley de Estabilidad Presupuestaria de 2002 limitó de forma consciente el estímulo del gasto público (pág. 42).
        \item \textbf{Deterioro de la Demanda Externa:} La absorción de recursos por parte de la demanda interna debilitó las exportaciones y disparó las importaciones, ampliando el déficit exterior debido a la inflación diferencial, un crecimiento nulo de la productividad y la pérdida de competitividad frente a países emergentes (pág. 42).
    \end{itemize}

    \item \textbf{Problemas estructurales del patrón de crecimiento 1995-2007:}
    \begin{itemize}
        \item Especialización productiva concentrada en sectores maduros de baja o media-baja tecnología que no generaban avances reales de productividad (pág. 43).
        \item Creación de un tejido empresarial atomizado, con serias debilidades organizativas y escasa proyección internacional (pág. 43).
        \item Profunda dualidad en el mercado de trabajo y un desajuste crónico entre salarios y productividad (pág. 43).
        \item Un elevado riesgo en el sistema financiero, que contrajo abundante deuda externa a corto plazo para financiar activos e hipotecas inmobiliarias a largo plazo (pág. 44).
        \item Ineficiencia y duplicidades en la gestión del sector público autonómico, inversión en infraestructuras improductivas y parálisis en reformas estructurales clave como las pensiones (pág. 44).
    \end{itemize}

    \item \textbf{Factores externos que propiciaron la crisis financiera de 2008-2014:}
    \begin{itemize}
        \item Una intensa y brusca escalada internacional en los precios del petróleo y de las materias primas (pág. 45).
        \item El estallido de la burbuja inmobiliaria en EE. UU., canalizado a través del colapso de las hipotecas \textit{subprime} y la ingeniería de títulos de alto riesgo mal calificados (pág. 45).
        \item El contagio de una crisis psicológica global de desconfianza que paralizó las decisiones de consumo e inversión a nivel mundial (pág. 45).
    \end{itemize}

    \item \textbf{Comportamiento de los desequilibrios internos en la crisis de 2008:}
    \begin{itemize}
        \item Abundancia de liquidez: \textbf{SÍ} (procedente del exterior, acumulada intensamente durante el ciclo de bonanza previo) (pág. 46).
        \item Subidas de los tipos de interés por parte del BCE al inicio de la crisis: \textbf{SÍ} (ejecutadas en 2008 con el fin de contener las presiones inflacionistas globales) (pág. 46).
        \item Interrupción de los flujos de capitales a España: \textbf{SÍ} (debido a la súbita pérdida de confianza de los inversores internacionales en la periferia europea) (pág. 47).
        \item Política monetaria restrictiva por parte del BCE: \textbf{NO} (al contrario, el BCE reaccionó aplicando una política monetaria netamente expansiva para proveer de liquidez de emergencia al sistema) (pág. 47).
        \item Aumento del crédito a empresas y familias: \textbf{NO} (se experimentó una severa y prolongada contracción o restricción crediticia) (pág. 47).
        \item Déficit presupuestario en las Administraciones Públicas: \textbf{SÍ} (la caída de los ingresos tributarios y el alza del gasto en estabilizadores automáticos abrieron una profunda brecha fiscal) (pág. 48).
        \item Gran aumento del desempleo: \textbf{SÍ} (ocasionado por la rigidez institucional y el ajuste centrado casi exclusivamente en el despido y la contratación temporal) (pág. 48).
        \item Aumento del déficit exterior: \textbf{NO} (el hundimiento de la demanda interna contrajo con fuerza las importaciones, reduciendo el déficit exterior hasta revertirlo en un saldo positivo desde 2012) (pág. 49).
    \end{itemize}

    \item \textbf{Factores acontecidos durante la recuperación económica 2014-2019:}
    \begin{itemize}
        \item Creación de empleo y subidas de salarios: \textbf{NO} (se registró una intensa creación de empleo, pero bajo un marco continuo de estricta moderación salarial) (pág. 51).
        \item Niveles de inflación moderados, aunque crecientes: \textbf{SÍ} (pág. 51).
        \item Aumento del nivel de endeudamiento de las familias y empresas no financieras: \textbf{NO} (los agentes económicos privados acometieron un intenso proceso de desapalancamiento financiero) (pág. 52).
        \item Déficit en el saldo de la balanza por cuenta corriente: \textbf{NO} (la balanza por cuenta corriente registró saldos positivos y superávits sostenidos debido a ganancias competitivas) (pág. 52, 53).
        \item Mejora en la capitalización de las entidades bancarias: \textbf{SÍ} (fruto directo de los procesos previos de reestructuración, saneamiento y recapitalización) (pág. 52).
        \item Reducción de la tasa de paro de larga duración: \textbf{NO} (se mantuvo en niveles estructuralmente altos, afectando a más del 50\% del total de desempleados) (pág. 53).
        \item Elevado nivel de endeudamiento exterior: \textbf{SÍ} (siguió operando como una de las principales debilidades de la economía) (pág. 53).
        \item Deuda pública inferior al 100\% PIB: \textbf{SÍ} (se situó en el 95,5\% en 2019, logrando bajar del umbral del 100\% aunque permaneciendo cercana a él) (pág. 50, 53).
        \item Crecimientos de productividad superiores a los de países desarrollados: \textbf{NO} (la productividad española continuó registrando tasas de variación inferiores a las de su entorno) (pág. 53).
    \end{itemize}

    \item \textbf{Medidas desarrolladas en España frente a la crisis sanitaria de la Covid-19:} \\
    El gobierno desplegó un paquete extraordinario de choque articulado bajo un triple objetivo macroeconómico:
    \begin{itemize}
        \item \textbf{Garantizar la actividad de las empresas:} Mediante la provisión masiva de líneas de avales públicos a través del Instituto de Crédito Oficial (ICO), seguros para firmas exportadoras y el aplazamiento de obligaciones fiscales y cotizaciones (pág. 61).
        \item \textbf{Mantener el empleo:} A través de la agilización institucional y simplificación de los Expedientes de Regulación Temporal de Empleo (ERTE), asumiendo el Estado la fuerza mayor para evitar despidos masivos (pág. 61).
        \item \textbf{Dar cobertura a la población desfavorecida:} Mediante la creación estructural del Ingreso Mínimo Vital (IMV) como red de seguridad de mínimos frente al impacto social de la parálisis económica (pág. 61).
    \end{itemize}

    \item \textbf{Respuesta de la UE ante la pandemia en comparación con la crisis de 2008:} \\
    La estrategia de respuesta comunitaria fue diametralmente opuesta en ambos episodios (pág. 62). Mientras que en la crisis financiera de 2008 la UE impuso una estricta doctrina de austeridad fiscal y recortes severos del gasto público, ante la Covid-19 desplegó un masivo e histórico escudo expansivo basado en la inyección de estímulos y liquidez (pág. 62). Esta respuesta anticíclica incluyó de forma destacada:
    \begin{itemize}
        \item 220.000 millones de euros en programas de ayuda directa del Banco Europeo de Inversiones (BEI) para sostener la liquidez de las PYMEs (pág. 62).
        \item 240.000 millones de euros movilizados por el Mecanismo Europeo de Estabilidad (MEDE) destinados en exclusiva a la adquisición de material médico e infraestructura sanitaria (pág. 62).
        \item La puesta en marcha por parte del BCE del Programa de Compras de Emergencia frente a la Pandemia (PEPP), dotado con 1,3 billones de euros para la adquisición masiva de bonos soberanos de los países más afectados (pág. 62).
        \item La histórica constitución del fondo solidario paneuropeo \textit{Next Generation EU}, inyectando un volumen de 750.000 millones de euros para financiar los planes nacionales de recuperación y transformación (pág. 62).
    \end{itemize}

\end{enumerate}
